\newpage
\section{Ponto neutro e margem estática}

O ponto neutro foi obtido do comando \texttt{st} do AVL na condição do
ponto de projeto, com a incidência de empenagem trimada, resultando em
$x_{np} = 30{,}125$ m, ou 74,1\,\%MAC. As margens estáticas
correspondentes às duas posições de CG são
\begin{equation*}
    SM_{\mathrm{fwd}} = \frac{x_{np} - x_{cg,\mathrm{fwd}}}{c_{\mathrm{ref}}}
    = 58{,}7\,\%\mathrm{MAC},
    \qquad
    SM_{\mathrm{aft}} = \frac{x_{np} - x_{cg,\mathrm{aft}}}{c_{\mathrm{ref}}}
    = 34{,}9\,\%\mathrm{MAC}.
\end{equation*}

Os valores são bem maiores que os estimados pelo \texttt{designTool} no
Lab 02, que situa o ponto neutro em 45,4\,\%MAC, com margens de 30,0\% e
6,2\%. A diferença é esperada e tem origem na fuselagem, que não é
modelada no AVL. O momento desestabilizante do corpo desloca o ponto
neutro real para a frente, e o \texttt{designTool} contabiliza esse efeito
por correlação empírica. Para o projeto, a estimativa do
\texttt{designTool} é a conservadora, e as margens do AVL devem ser lidas
como as da combinação asa e empenagem isoladas. Essa mesma diferença
explica as incidências de trimagem elevadas da Seção 4 e sugere, como
próximo passo, incluir a fuselagem no modelo do AVL para as análises
finais de estabilidade.
